% =============================================================================
% sec_prediction.tex -- Pre-Engagement Prediction  (ToG manuscript, item M2-prediction, 2026-09-14)
%
% Answers CoG 2026 reviews: R3 (motivate every method for a non-expert: task, grouped cross-validation,
% AUC, match bootstrap, attribution forensics), R1 (results must not hinge on hand-set constants: the G x D
% sweep and the presence-gate points, both filled 2026-09-14 by item F3-fill-results, sec:pred-constants) and the audit-driven withdrawal that the
% Introduction announces (reviewer_response_matrix.md rows R1-2, R3-5).
%
% FACTS THIS FILE MUST KEEP (checked against code and artefacts by this item):
%   * The CoG submission's text has no scale classes (REPO/paper/abstract.tex, intro.tex, methods.tex,
%     result.tex: no class definition, no n_min, no per-class result; the words occur only descriptively:
%     abstract.tex l.8 "teamfights", intro.tex l.5 "skirmishes or teamfights", methods.tex l.60 "isolated
%     picks"; re-read 2026-09-14).  The scale conclusion came from the post-conference
%     extension draft REPO/paper/tog/main_en.tex (v0.2, SUPERSEDED banner, lines 2-8, 20, 30, 42, 52, 60-62)
%     and from re-runs on the corpus-v3 shards (REPO/docs/DATA_MANIFEST.md "Corpus v3").  Never write that
%     the conference paper claimed a scale gradient.
%   * What IS a CoG-era issue: the conference pipeline used the legacy feature path.  REPO/core/presets.py
%     PRESETS["cog2026"] (tag v1.0-cog2026): TIME_NORM_ABSOLUTE False, ANCHORS_CAUSAL False,
%     TAB_FRAME_AGE_FEATURE False; PRESETS["v3.3"]: True, True (TIME_NORM_DENOM_MIN 45.0), True.
%     git show v1.0-cog2026:gameplay/pipeline_cache.py line 48: time_norm = float(t / max(1.0, T - 1))
%     (re-read 2026-09-14).  Fix commit 5b5f9c8 (2026-09-09 09:33, "Input feature audit: two leaks fixed,
%     spatial block repaired, frame age added").
%   * "No scale gradient" is never written as absolute (A6 cut table).
%
% Provenance.  Every number carries a comment with its artefact and key; derived numbers show the
% arithmetic.  All values were re-read from the artefacts on 2026-09-14 by this item.  Path prefixes:
%   FEAT/    = D:/LOL_Project/fusion_2615/features/
%   REV/     = D:/LOL_Project/fusion_2615/features/tog_revision/
%   REPO/    = C:/Users/todtj/PycharmProjects/LOL_teamfight/
%   WT/      = C:/Users/todtj/문서/LOL_Teamfight/worktrees/engagement-state-value/
% No comment cites a session scratchpad (not archived, not durable).  Every printed value is either in an
%   archived artefact under FEAT/ or REV/, or recomputable from committed code and archived inputs, with the
%   computation given where it is used (the one case: one model's pick-minus-teamfight interval on the patch
%   holdout, Section sec:pred-scoreboard).
% Protocols are never mixed in one table or one difference:
%   "five-fold"      = match-grouped 5-fold out-of-fold predictions on the whole v3.3 corpus, published
%                      LightGBM configuration (FEAT/scale_decomposition_v33_market_event.json and .preds.npz;
%                      REV/A6-definition-sensitivity/);
%   "patch holdout"  = train 15.14 / early stopping 15.15 / test 15.16 (FEAT/model_comparison_v33_patch.json
%                      and .preds.npz; REV/A8-input-audit-and-cis/; REV/A10-perclass-contrast/).
% Learner of the five-fold predictions: WT/scripts/run_scale_decomposition.py oof_predictions():
%   GroupKFold(n_splits=5) on match ids (no shuffle), LGBMClassifier(n_estimators=400, learning_rate=0.05,
%   num_leaves=31, subsample=0.9, colsample_bytree=0.9, random_state=7), no early stopping.  subsample_freq is
%   left at its default 0, and LightGBM 4.6 enables row subsampling only when subsample_freq > 0
%   ("subsample_freq ... <=0 means no enable", LGBMClassifier docstring, anaconda python used for the runs),
%   so every tree sees all training rows and only the 90 % column draw is active.  The text says so.
% Match bootstrap of the five-fold file: same script, cluster_bootstrap(): rng default_rng(7), resample
%   match ids with replacement, roc_auc_score on the concatenated rows, 2.5 / 97.5 percentiles, n_boot 1000.
%
% Cross-references.  \ref for labels in committed files (sec:intro, sec:background, sec:definition,
%   sec:def-constants, sec:def-killless, sec:limitations, subsec:limitations, sec:availability);
%   \fwdref (macros.tex) for sec:label and for the labels of sec_learners.tex.  sec_label.tex now exists
%   (\label{sec:label} at l.26) but is untracked, and sec_learners.tex is untracked and partial (sec:learners,
%   sec:learners-input, sec:learners-ladder, sec:learners-protocol, sec:learners-terms); \fwdref resolves to
%   \ref in a full build and prints a placeholder only if a file is missing.  sec:limit is not referenced.
% Labels defined here: sec:prediction, sec:pred-task, sec:pred-protocol, sec:pred-headline, sec:pred-audit,
%   sec:pred-cut, sec:pred-scoreboard, sec:pred-constants; tab:pred-headline, tab:pred-leak, tab:pred-cut,
%   tab:pred-scoreboard, tab:pred-gxd (added by item F3-fill-results).
% Build check 2026-09-14 (after the fact-check fixes): compiled with IEEEtran + macros.tex + bibtex in the
%   scratchpad (stub sections for the \ref / \fwdref targets): no errors, no overfull boxes, no undefined
%   references, all three citations resolve.
% Fact-check round (2026-09-14): cut qualifier and cut-3/5 full-model gaps in sec:pred-scoreboard; "-1
%   stored count" wording in three captions; 0.675 source; spatial inventory; interpolation of positions;
%   tab:pred-leak caption; durable provenance of one-model class gaps; plain-words LightGBM and early
%   stopping (R3); "makes no claim" wording.  No \pending marker could be filled (see the four keys below).
% \pending keys used here (underscores escaped for text mode; grep "pending" finds them):
%   leak_ablation_ci.  shap_forensics was filled on
%   2026-09-14 by item F2-fill-and-gloss from REV/A7-shap-forensics/rerun.json (see sec:pred-audit); leak_ablation_ci stays
%   pending (REV/C3-leak-ablation holds only smoke_m300 on 2026-09-14 16:10).  sec:limit-shap is now referenced.
%   gxd_sweep and presence_gate_v33 were filled on 2026-09-14 by item F3-fill-results from
%   REV/A6-definition-sensitivity/gd_sweep/gd_sweep_summary.json and REV/A6-definition-sensitivity/presence_gate/
%   presence_gate_summary.json (sec:pred-constants; new label tab:pred-gxd).
% Citation keys (all in REPO/docs/references/definition_refs.bib, marked VERIFIED): efron1993bootstrap,
%   ke2017lightgbm, lundberg2020local.
% =============================================================================

\section{Pre-Engagement Prediction}\label{sec:prediction}

This section states the prediction task, the protocol that evaluates it and the result. It then reports
the audit of the model inputs, which withdrew an earlier conclusion that predictability depends on
engagement scale; shows that any comparison between scale classes depends on where the class boundary
is placed; and measures how much of the result a model that reads only the scoreboard reproduces.
% Gloss at first use (item F2, 2026-09-14): the full description is in sec:pred-protocol ("The learner is LightGBM").
Learners other than gradient-boosted trees, many small decision trees added one after another with each
fitted to the errors of those before it (Section~\ref{sec:pred-protocol}), are compared in
Section~\fwdref{sec:learners}, and the label
and its variants are examined in Section~\fwdref{sec:label}.

% -----------------------------------------------------------------------------
\subsection{The Task}\label{sec:pred-task}

In plain words, the task is to state, shortly before an engagement starts, how likely the blue team is
to win it, using only what the match telemetry has recorded up to that moment.

% B = 15 s: FEAT/model_comparison_v33_patch.json corpus_manifest.detector.TF2_ENGAGE_PRE_KILL_MS 15000;
%   the presence gate is evaluated at the same tau (sec_definition.tex lines 228-231).
% 33,905 of 566,452: FEAT/model_comparison_input_audit.json truth.rows_total 566452, rows_labelled 532547;
%   566452 - 532547 = 33905.
% 532,547 / 191,940 / 50.8 %: FEAT/scale_decomposition_v33_market_event.json n 532547, n_matches 191940,
%   positive_rate 0.507854.
Formally, let $e$ be an engagement (Section~\ref{sec:definition}) whose first kill occurs at time
$t_1(e)$. Its \emph{cutoff} is $\tau(e) = t_1(e) - B$ with $B = 15$\,s, the instant at which the
definition's presence gate is evaluated. The target $y(e)$ is 1 when the blue team wins the engagement
under the \texttt{market\_event} label (Section~\fwdref{sec:label}) and 0 when the red team wins it;
engagements that the label scores as a draw have no target and are excluded, 33,905 of 566,452 detected
engagements. A model maps an input vector $x(e)$, computed only from telemetry recorded up to $\tau(e)$,
to a number $\hat p(e) \in [0, 1]$, its estimate of $\Pr\{y(e) = 1 \mid x(e)\}$. The corpus holds
532,547 labelled engagements from 191,940 matches, and the blue team won 50.8\% of them.

% Window and bins: REPO/core/config.py FIGHT_CONTEXT_SEC 30 (l.432), BIN_MS 5000 (l.321); 30 / 5 = 6 bins.
% Base-quantity groups: REPO/docs/INPUT_FEATURE_AUDIT_V3.md section 1 (inventory, 2026-09-09):
%   per player slot 17 snapshot + 11 state + 11 rune + 25 champion statistic + 12 damage + 16 item = 92;
%   92 x 10 = 920; game-wide time_norm 1 + bans 10 + team differences 15 + event counts 44 + spatial
%   (6 + 5 + 14 =) 25 = 95; 920 + 95 = 1015 (= FEAT/deep_tabular_v33_patch_full.json n_bases 1015).
% Spatial 25 (REPO/gameplay/feature_spatial.py SPATIAL_FEATURE_NAMES l.17-43; audit section 1 rows l.22-24):
%   6 = pos_fight/pos_blue/pos_red x,y, zeroed by [P0-1] (REPO/gameplay/features.py l.23-30 and l.564-571,
%   ZERO_XY_IN_EXTRA_SEQ True at REPO/core/config.py l.509; the zeroed copy enters macro_seq, which the
%   "full" and "tri_modal" feature sets read, features.py l.573 and l.613-623; audit l.22 "fixed at 0 ->
%   constant, removed");
%   5 = dist_team_sep, standoff_pairs_frac, mean_min_enemy_dist, d_standoff_pairs_frac, d_mean_min_enemy_dist;
%   14 = dist_obj_nearest, near_obj x5, dist_tower_nearest, in_tower_range, near_tower_radius, zone x5.
% Frames before tau only: REPO/gameplay/pipeline.py l.161-175 node_max_snapshot_ms = label_start_ms - 1
%   (label_start_ms = engage_ts = tau, l.128) passed to REPO/gameplay/pipeline_interp.py
%   interpolate_node_global / interpolate_abs_xy, which cap the query time at the last frame with timestamp
%   <= tau - 1 ms ("Prevent future-frame interpolation leakage", l.117-122); scalars ffill
%   (INTERP_SCALARS_METHOD "ffill", REPO/core/config.py l.337); coordinates linear between the two frames
%   around the (capped) query, alpha clipped to [0, 1] (l.127-138; _interp_xy_guarded holds on a jump
%   > 7,000 u or a death, audit section 1 row 1).  The x_seq of feature set "full" is node_flat + macro_seq
%   (features.py l.613-620), i.e. 920 + 95 base quantities.  Frames every ~60 s: sec_background.tex (60.00-60.17 s in the corpus).
% Seven statistics: REPO/docs/FEATURES.md l.279-297; FEAT/model_comparison_input_audit.json check "suffix
%   families are contiguous and suffix-major" ['last','mean','std','min','max','delta','slope'].
% 7,106 / 942 / 6,164: FEAT/model_comparison_input_audit.json truth.columns_total 7106, columns_constant 942,
%   columns_non_constant 6164; 1015 * 7 + 1 = 7106.  The tree pipeline's filter reads column ranges over all
%   shard rows and never the labels (sec_learners.tex, paragraph "Columns").
The input summarises the \emph{observation window} $[\tau - 30\,\mathrm{s},\ \tau]$, divided into six
bins of 5\,s. Of its 1,015 base quantities, 920 describe the ten players (for each: 17 snapshot quantities
such as position, level, experience, gold and creep score; 11 buff and ability states; 11 rune slots; 25
champion statistics; 12 cumulative damage totals; 16 item quantities); the other 95 are game time, 10
champion bans, 15 blue-minus-red team differences, 44 counts of timeline events and 25 spatial
quantities. Of the spatial quantities, five measure how far apart the two teams stand and 14 where the
fight lies relative to towers, objectives and map zones; the other six are coordinates of the fight and
team centres, which the input holds at zero, so they are among the constant columns counted below.
Quantities read from participant frames use only frames recorded before $\tau$. Frames arrive about every
60\,s, so a quantity carried forward from the
last frame changes at most once inside the window. Champion positions are instead interpolated linearly
between consecutive frames, but never past the last frame before $\tau$: they can change in every bin up
to that frame and keep its value after it. Events carry millisecond timestamps and are counted in the bin
that contains them.
Each base quantity is summarised over the six bins by seven statistics: the last value, the mean, the
standard deviation, the minimum, the maximum, the change from the first to the last bin, and the
least-squares slope. A final column, \emph{frame age}, gives the time from the last frame to $\tau$ and so
tells a model how stale the frame quantities are. Each engagement is thus a row of
$1{,}015 \times 7 + 1 = 7{,}106$ numbers. Of these columns, 942 take a single value in every engagement of
the corpus and carry no information; the tree learners use the other 6,164
(Section~\fwdref{sec:learners-input}).

% "who wins an imminent engagement, not whether one occurs": sec_background.tex term "Prediction cutoff".
\emph{What a leak is here.} An input \emph{leaks} when its value at $\tau$ depends on anything recorded
after $\tau$: the engagement itself, later events, or properties of the whole match such as its length or
its winner. A leaking input answers part of the question from the future. The AUC it produces cannot be
reached by a model used before the fight, and a comparison built on it, between scale classes, learners or
labels, may measure how much each group gains from the leak rather than how predictable the group is.
Leaks need not come from a column that obviously reads the future; a scale factor or a reference table
computed from the complete match leaks just as well, and Section~\ref{sec:pred-audit} reports two of that
kind. One use of the future is part of the task and is not a leak in this sense: $\tau$ is placed from a
kill that has not yet happened, so the task is to predict who wins an engagement known to be imminent,
not whether one will occur (Section~\ref{subsec:limitations}).

% -----------------------------------------------------------------------------
\subsection{Evaluation Protocol}\label{sec:pred-protocol}

\emph{Grouped cross-validation.} A model's skill has to be measured on engagements it was not trained on.
Five-fold cross-validation divides the data into five parts, the \emph{folds}; a model trained on four
folds predicts the engagements of the fifth, and repeating this for every fold gives each engagement one
\emph{out-of-fold} prediction from a model that never saw it. The folds are formed from whole matches, so
that all engagements of a match fall in the same fold. Engagements of one match share players, champions
and much of the game state (a gold lead persists from one fight to the next), and a model trained on some
engagements of a match could partly recognise the match when predicting its other engagements rather than
predict the fight. The AUC is computed once over all 532,547 out-of-fold predictions.
% Learner: see header (oof_predictions).  Configuration text shared with sec_learners.tex paragraph
%   "Gradient-boosted trees"; the row subsample inactive per the header note.
The learner is LightGBM~\cite{ke2017lightgbm}, an implementation of gradient-boosted decision trees. A
decision tree predicts by a short sequence of yes-or-no questions about the input columns; boosting builds
many small trees one after another, fits each new tree to the errors that the trees before it still make,
and adds up their outputs, scaled by a learning rate, to form the prediction. It is used in its published
configuration (Section~\fwdref{sec:learners-ladder}): 400 trees of at most 31 leaves, learning rate 0.05,
and a random 90\% of the columns drawn for each tree. The configuration also sets a row fraction of 90\%, but LightGBM
applies it only when a subsampling frequency is set, which this configuration leaves at its default, so
every tree sees all training rows. Each fold's model is fitted once, with one random seed.

% Patch holdout: FEAT/model_comparison_v33_patch.json split {train 15.14, val 15.15, test 15.16},
%   rows {187547, 191184, 153816}.  lgbm_paper test AUC 0.666484 (models.lgbm_paper.overall_auc).
% Early stopping: WT/scripts/run_model_comparison_v33.py run_model() l.130-131: fit(eval_set=[15.15],
%   eval_metric="auc", callbacks=[early_stopping(100)]); with first_metric_only left False the callback
%   watches every validation metric (AUC and LightGBM's default binary log-loss) and stops when one has not
%   improved for 100 trees; best_iteration_ is used for prediction.  Hence "score", not "AUC", in the text;
%   plain-words definition shared with sec_learners.tex "Terms".
% Training-set sizes: five-fold 532547 * 4 / 5 = 426,038 (about 426,000); patch 187,547.
\emph{Why this protocol for the headline, and when the other.} Cross-validation scores every engagement,
which the class analyses below need, and it measures predictability within the three patches of the
corpus. It does not measure forecasting across a patch change, because every fold mixes all three
patches. The \emph{patch holdout} does: models are trained on patch 15.14 and tested once on 15.16
(Section~\fwdref{sec:learners-protocol}), and patch 15.15 is used for \emph{early stopping}: trees stop
being added once the model's score on patch 15.15 has not improved for 100 consecutive trees, and the
model at its best point is kept; more trees would fit accidental regularities of the training patch
rather than structure that carries over. The patch holdout is used for the learner comparison and
for the scoreboard contrast of Section~\ref{sec:pred-scoreboard}. The published configuration scores
0.6699 under cross-validation and 0.6665 under the patch holdout, where it also stops early. The
difference mixes the size of the training set (about 426,000 against 187,547 engagements), early stopping
and the patch change, so it is not interpreted, and the two protocols never share a table or a
difference.

% AUC 0.6699 -> "about two times in three": 0.6699 ~ 2/3 = 0.6667.
\emph{Area under the ROC curve (AUC).} The model outputs a probability, and the AUC evaluates how well
those probabilities rank engagements: it is the probability that a randomly chosen engagement that blue
won receives a higher predicted probability than a randomly chosen engagement that blue lost, with ties
counted as one half. It equals 0.5 for a model that ranks at random and 1 for a perfect ranking,
whatever the share of blue wins, and it needs no decision threshold. An AUC of 0.67 means that the model
orders such a pair correctly about two times in three. The AUC does not measure calibration, that is,
whether engagements given a probability of 0.7 are won 70\% of the time.

% Replicate counts: FEAT/scale_decomposition_v33_market_event.json bootstrap.*.n_boot 1000;
%   REV/A6-definition-sensitivity/scale_cut_sensitivity_v33.json provenance.n_boot 400, seed 7;
%   REV/A10-perclass-contrast/perclass_lead_contrast_v33.json (A8 method, seed 7, 2,000 replicates).
% "not refitted", "unadjusted": REV/A8-input-audit-and-cis/paired_learner_cis_v33.json deviations
%   ("each learner is one fitted model ... intervals cover test-set sampling only"; "intervals are per
%   comparison and unadjusted"); REV/A10 .../perclass_lead_contrast_v33.json deviations (same wording).
% Matches resampled: 191,940 = FEAT/scale_decomposition_v33_market_event.json n_matches; 55,449 =
%   REV/A10-perclass-contrast/perclass_lead_contrast_v33.json overall.n_matches (= A8 comparisons.*.n_matches).
\emph{Match bootstrap.} Another sample of matches from the same population would give a somewhat
different AUC, and a confidence interval states how different. The bootstrap estimates this variation
from the data at hand~\cite{efron1993bootstrap}. It draws, with replacement, as many matches as are
evaluated (all 191,940 under cross-validation, the 55,449 of the test patch under the patch holdout),
keeps every engagement of each drawn match as often as the match was drawn, and recomputes the AUC from
the stored predictions; after many such resamples, the 2.5th and 97.5th percentiles of the recomputed
AUCs form the 95\% interval. Matches, not engagements, are resampled because engagements of one match are
not independent, and resampling them one by one would count correlated engagements as
separate evidence and make the intervals too narrow. For a difference between two AUCs on the same
engagements, such as two classes or two models, both are computed on every resample and the interval is
taken over the resampled differences; such an interval is called \emph{paired}. Models are not refitted
inside the resamples, so the intervals describe the sampling of the evaluated matches, not the variation
from retraining. Intervals use 1,000 resamples for the classes in Table~\ref{tab:pred-headline}, 400 for
its first row and for Table~\ref{tab:pred-cut}, and 2,000 for Table~\ref{tab:pred-scoreboard} and the
comparisons that accompany it, and they are not adjusted for the number of comparisons.

% -----------------------------------------------------------------------------
\subsection{Result}\label{sec:pred-headline}

% Overall: REV/A6-definition-sensitivity/scale_cut_sensitivity_v33.json overall.auc point 0.669909,
%   ci_2.5 0.668356, ci_97.5 0.671160 (n_boot 400, seed 7) -> 0.6699 [0.6684, 0.6712].  Point equals
%   FEAT/scale_decomposition_v33_market_event.json overall_auc 0.6699089924634439 (that file has no overall
%   interval).
% Classes: FEAT/scale_decomposition_v33_market_event.json by_participation_scale.{pick,skirmish,teamfight}
%   n 101798 / 320878 / 109829, share 0.191153 / 0.602535 / 0.206233, auc 0.678869 / 0.663327 / 0.680888;
%   bootstrap.{pick,skirmish,teamfight} ci_2.5 / ci_97.5 0.675686 / 0.682056, 0.661409 / 0.665259,
%   0.677646 / 0.684035 (n_boot 1000).  The 42 rows whose stored participation count is -1 are in no class
%   (532547 - 101798 - 320878 - 109829 = 42; 42 / 532547 = 0.0079 %).  -1 is how the index stores a count
%   of zero: REPO/gameplay/fights.py l.1465-1466 always writes an int count, and REPO/data/index_split.py
%   l.361-362 stores int(fight.get("det_cluster_blue", -1) or -1), so 0 becomes -1.  REV/A6 reads these
%   rows as zero participants (scale_participation_v33.json definitions.negative_counts); REPO/scripts/
%   run_scale_decomposition.py scale_class reads them as unknown (scale_cut_sensitivity_v33.json deviations[3]).
% Cut: FEAT/scale_decomposition_v33_market_event.json teamfight_min 4.
With the protocol above, LightGBM ranks engagement outcomes at an AUC of 0.6699 (95\% interval 0.6684 to
0.6712; Table~\ref{tab:pred-headline}). Table~\ref{tab:pred-headline} also scores the same out-of-fold
predictions within the scale classes of Section~\ref{sec:definition}: picks, in which the smaller side has
at most one participating champion ($n_{\min} \le 1$), skirmishes ($2 \le n_{\min} \le 3$) and teamfights
($n_{\min} \ge 4$). Scale is known only after the engagement, so it is never an input, and no model is
trained per class. The class AUCs are 0.6789, 0.6633 and 0.6809. Whether and how they differ depends on
where the teamfight class begins (Section~\ref{sec:pred-cut}), and an earlier analysis that ranked the
classes was built on leaking inputs (Section~\ref{sec:pred-audit}).

\begin{table}[t]
\centering
\caption{Pre-engagement prediction, five-fold cross-validation grouped by match: out-of-fold AUC of the
published LightGBM configuration, overall and within scale classes defined on $n_{\min}$, the smaller
side's number of participating champions. One
set of predictions is scored on every row. Intervals: match bootstrap, 400 resamples (overall) and 1,000
(classes). The 42 engagements whose stored participation count is $-1$ (the index stores a side with no
participating champion as $-1$) are in no class.}
\label{tab:pred-headline}
% Rows: see the comments above the paragraph.  Shares: 101798 / 532547 = 19.12 %, 320878 / 532547 = 60.25 %,
%   109829 / 532547 = 20.62 %.
\footnotesize
\setlength{\tabcolsep}{4pt}
\begin{tabular}{@{}lrrc@{}}
\toprule
Engagements & $n$ & Share (\%) & AUC [95\% CI] \\
\midrule
All & 532,547 & 100.00 & 0.6699 [0.6684, 0.6712] \\
Pick ($n_{\min} \le 1$) & 101,798 & 19.12 & 0.6789 [0.6757, 0.6821] \\
Skirmish ($2$--$3$) & 320,878 & 60.25 & 0.6633 [0.6614, 0.6653] \\
Teamfight ($\ge 4$) & 109,829 & 20.62 & 0.6809 [0.6776, 0.6840] \\
\bottomrule
\end{tabular}
\end{table}

% -----------------------------------------------------------------------------
\subsection{The Input Audit and a Withdrawn Conclusion}\label{sec:pred-audit}

% Audit: REPO/docs/INPUT_FEATURE_AUDIT_V3.md sections 1-2 (2026-09-09): table of every base-quantity group
%   with source, time alignment, fill and verdict; "누수 2건 수정, 무의미 블록 1건 수정, 특징 1개 추가"
%   (two leaks fixed, one meaningless block repaired, one feature added); commit 5b5f9c8.
% Leak 1: cache time_norm = t / (T - 1) (git show v1.0-cog2026:gameplay/pipeline_cache.py l.48); REPO/core/
%   config.py l.435-440 ("normalised by the match's own length, which is not known at the cutoff"; with
%   TIME_NORM_ABSOLUTE the feature path rewrites it as min(t / 45 min, 1) and the phase encoding recovers
%   minutes with the same constant).
% Leak 2: REPO/core/config.py l.441-445 (meta["anchors"] = "every tower that fell and every objective killed
%   anywhere in the match"; ANCHORS_CAUSAL uses the static map and only buildings destroyed at or before the
%   cutoff, REPO/gameplay/anchors.py causal_anchors / standing_towers).
% Non-leak repair not toggled below: spatial block computed from centre-relative coordinates, repaired by an
%   absolute-coordinate path that is unconditional (REPO/gameplay/pipeline.py l.159-276 xy_abs_seq; no flag).
% 0.675: REPO/paper/abstract.tex l.18 ("AUC~$0.675$"); REPO/paper/result.tex l.16 ("(AUC~.675)") and l.31
%   (Table tab:main-results, "LightGBM (engineered tabular) & .675", held-out patch 15.16).
To make sure that the result above uses no future information, every one of the 1,015 base quantities was
checked for where it comes from, which instant it describes, and how it is filled between frames. The
audit found two leaks. First, game time entered the input divided by the length of the whole match, a
number that is not known at the cutoff; the quantities derived from it inherited the leak. Second, the
map quantities that measure how far a fight lies from towers and objectives were computed against every
tower destroyed and every objective taken anywhere in the match, so a distance to the nearest fallen tower
could reflect a tower that fell after the cutoff. Both were fixed: game time is now elapsed minutes
divided by 45 and capped at 1, and the map reference is the static map with only the buildings destroyed
at or before the cutoff. The audit also repaired map quantities that had been computed from coordinates
relative to the fight centre, which is not a leak, and added the frame-age column. The conference pipeline
used both leaking inputs (the code keeps its settings as a named configuration preset), so its reported
AUC of 0.675 is not a clean reference, and it is not compared with the results here, which also differ
from it in definition, label and protocol.

% FEAT/leak_ablation_v33.json (script REPO/scripts/run_leak_ablation.py, docstring; run 2026-09-09, log
%   FEAT/leak_ablation_v33.log): n_matches 5000 (random.Random(7).sample of cached matches), n_refs 13422,
%   every config n_rows 12603, n_features 7106; rows asserted identical across configs
%   ("row set changed between configurations"); label market_event with the v3.3 detector constants
%   (V3 dict: G 13,700 ms, D 4,264 u, R 1,600 u, B 15 s, H 35 s, attribution "engagement", ties dropped);
%   learner sd.oof_predictions = the published configuration, GroupKFold(5) on match ids; teamfight =
%   smaller side >= 4.  Fold training sets 10,082-10,083 engagements (log "fold k/5: train=...").
% Class sizes: configs.clean.by_class pick n 2419, skirmish 7556, teamfight 2628.
\emph{How much the leaks were worth.} A controlled ablation measures it. On 5,000 randomly drawn matches,
12,603 engagements were built four times under the current definition and label, with each leaking input
switched on or off and everything else fixed, and each build was scored by the protocol of
Section~\ref{sec:pred-protocol}. The engagements, labels and folds are identical across the four builds,
so the differences are due to the two inputs alone. Each fold trains on about 10,000 engagements, against
about 426,000 in Table~\ref{tab:pred-headline}, and the absolute AUCs are lower; only the differences
between builds are read.

% Differences (configs.X minus configs.clean; computed 2026-09-14 from the JSON):
%   time_leak:   overall 0.639154 - 0.619726 = +0.019428; pick 0.620699 - 0.618075 = +0.002624;
%                skirmish 0.626096 - 0.612278 = +0.013818; teamfight 0.691158 - 0.642066 = +0.049092;
%                pick_minus_teamfight -0.070459 vs clean -0.023991: change -0.046468.
%   anchor_leak: overall +0.004375; pick +0.002525; skirmish +0.002024; teamfight 0.653579 - 0.642066 =
%                +0.011513; pick_minus_teamfight -0.032978: change -0.008988.
%   both_leaks:  overall 0.639965 - 0.619726 = +0.020239; pick 0.623716 - 0.618075 = +0.005641;
%                skirmish 0.627945 - 0.612278 = +0.015667; teamfight 0.686937 - 0.642066 = +0.044871;
%                pick_minus_teamfight -0.063221: change -0.039230.
Table~\ref{tab:pred-leak} gives the result. Dividing game time by match length raised the overall AUC by
0.019, but not evenly: it raised the AUC of teamfights by 0.049, of skirmishes by 0.014 and of picks by
0.003, and so widened the difference between the pick and teamfight AUCs from $-0.024$ to $-0.070$. The
map anchors alone raised the teamfight AUC by 0.012 and the pick and skirmish AUCs by at most 0.003. With both leaking
inputs, as in the conference pipeline, the overall AUC rose by 0.020, the teamfight AUC by 0.045 and the
pick AUC by 0.006, and the pick-minus-teamfight difference reached $-0.063$. A leak can therefore create a
difference between classes as well as inflate a number. We have not tested why teamfights gain most;
they also occur later in the match (median cutoff at minute 21.1, against 8.7 for picks), so the class
and the game time at which the leaking time input is read are confounded.
% Median cutoff minutes: REV/A6-definition-sensitivity/scale_participation_v33.json
%   populations.labelled.participation.cuts."4".teamfight.median_engage_minute 21.148,
%   .pick.median_engage_minute 8.718 (engage_ts = cutoff tau: REPO/scripts/build_corpus_shard.py l.66
%   engage_ts=int(r.t_start_ts); REPO/gameplay/fights.py l.1371 engage_ts_val = first_kill_ts - pre-kill).
%   Whole labelled corpus, cut 4.  The explanation is flagged as untested in the text.
% No interval: FEAT/leak_ablation_v33.json has no predictions and no bootstrap block; the script writes
%   only AUCs (run_leak_ablation.py main()).
The ablation kept no predictions, so these differences carry no interval
(\pending{leak\_ablation\_ci}{paired match-bootstrap intervals for the differences in
Table~\ref{tab:pred-leak}; needs scripts/run\_leak\_ablation.py re-run with its predictions saved}). With
2,419 picks and 2,628 teamfights, class AUCs in this sample are far less precise than in
Table~\ref{tab:pred-headline}, which is why the clean build's own difference of $-0.024$ is not read as a
finding; on the full corpus the same inputs give $-0.0019$ at this cut (Table~\ref{tab:pred-cut}).

\begin{table}[t]
\centering
\caption{Input-leak ablation on 12,603 engagements from 5,000 matches: the same engagements, labels and
match-grouped folds, with the two leaking inputs switched on or off. Time leak: game time divided by match
length. Anchor leak: map anchors from the whole match. Both leaks: both leaking inputs switched on, as in
the conference pipeline.
Out-of-fold AUC of the published LightGBM configuration; classes at $n_{\min} \ge 4$ ($n$ = 2,419 picks,
7,556 skirmishes, 2,628 teamfights). No intervals were computed.}
\label{tab:pred-leak}
% FEAT/leak_ablation_v33.json configs.{clean,time_leak,anchor_leak,both_leaks}: auc; by_class.pick.auc;
%   by_class.skirmish.auc; by_class.teamfight.auc; pick_minus_teamfight.
%   clean       0.619726  0.618075  0.612278  0.642066  -0.023991
%   time_leak   0.639154  0.620699  0.626096  0.691158  -0.070459
%   anchor_leak 0.624101  0.620601  0.614302  0.653579  -0.032978
%   both_leaks  0.639965  0.623716  0.627945  0.686937  -0.063221
% "legacy path" = both flags False = REPO/core/presets.py PRESETS["cog2026"] feature flags.
% both_leaks is NOT the conference pipeline's input: it is the v3.3 representation with the two flags off
%   (REPO/scripts/run_leak_ablation.py docstring "both_leaks ... (v3 feature path)", CONFIGS l.37; 7,106
%   columns incl. frame age, FEAT/leak_ablation_v33.json configs.both_leaks.n_features 7106; absolute-
%   coordinate spatial block; v3.3 definition and label).  The conference pipeline had 7,105 columns, no
%   frame age (PRESETS["cog2026"] TAB_FRAME_AGE_FEATURE False) and a centre-relative spatial block.  Hence
%   "as in the conference pipeline" qualifies only the two flags.
\footnotesize
\setlength{\tabcolsep}{3pt}
\begin{tabular}{@{}lccccc@{}}
\toprule
Inputs & All & Pick & Skirmish & Teamfight & Pick $-$ team. \\
\midrule
Audit-fixed & 0.6197 & 0.6181 & 0.6123 & 0.6421 & $-0.0240$ \\
Time leak only & 0.6392 & 0.6207 & 0.6261 & 0.6912 & $-0.0705$ \\
Anchor leak only & 0.6241 & 0.6206 & 0.6143 & 0.6536 & $-0.0330$ \\
Both leaks & 0.6400 & 0.6237 & 0.6279 & 0.6869 & $-0.0632$ \\
\bottomrule
\end{tabular}
\end{table}

% Withdrawn draft: REPO/paper/tog/main_en.tex (v0.2; SUPERSEDED banner l.2-8; title l.20 "... Decomposes by
%   Scale"; l.30 "pre-engagement predictability decomposes by scale"; l.42 teamfight n_min >= 3;
%   l.52 "teamfights reach 0.7842 against 0.7165 for picks ... gap of -0.0677 (95% CI [-0.0702,-0.0650])",
%   under market_lex; l.32 948,369 engagements).  Its corpus: REV/A6-definition-sensitivity/
%   scale_participation_v33.json v2_reference n 948369, "v2 corpus (G 18 s, D 4,000 u, R 1,800 u, B 10 s)",
%   i.e. the conference constants.
% Corpus-v3 re-run: FEAT/scale_decomposition_v3_market_event.json (file time 2026-09-09 08:08, before
%   5b5f9c8 at 09:33): n 535444, teamfight_min 4, y_key y_market_event, by_participation_scale pick auc
%   0.672614, teamfight auc 0.770143; 0.672614 - 0.770143 = -0.097529; bootstrap.pick_minus_teamfight
%   ci_2.5 -0.101877, ci_97.5 -0.093388 (n_boot 1000).  Corpus v3 = v3.3 constants (REPO/docs/DATA_MANIFEST.md
%   "Corpus v3": G 13.7 s, D 4,264 u, R 1,600 u, B 15 s) but legacy inputs and before the later label-window,
%   attribution and representation fixes (DATA_MANIFEST.md "Corpus v3.3"; REPO/docs/INPUT_FEATURE_AUDIT_V3.md
%   section 4).
% CoG text has no scale classes: REPO/paper/{abstract,intro,methods,result}.tex (grep: no pick, skirmish,
%   teamfight or n_min; result.tex stratifies by game phase and gold imbalance only).
\emph{The withdrawn conclusion.} The conference version reported no scale classes. After it, a draft of
this extension concluded that pre-engagement predictability decomposes by scale: with teamfights taken as
$n_{\min} \ge 3$, under a different label and on engagements detected with the conference constants,
teamfights reached an AUC of 0.7842 against 0.7165 for picks, a difference of $-0.0677$ (95\% interval
$-0.0702$ to $-0.0650$). A re-run on engagements detected with the current constants, but still with the
leaking inputs and before later fixes to the label and the input representation, gave 0.7701 for
teamfights ($n_{\min} \ge 4$) against 0.6726 for picks, a difference of $-0.0975$. Both analyses used the
leaking inputs. The ablation shows that, on identical engagements, those inputs widen the difference in
that direction, by 0.039 with both and by 0.046 with the time input alone, and on audit-fixed inputs the
difference at the adopted cut is $-0.0019$
(95\% interval $-0.0065$ to $+0.0028$). We therefore withdraw the conclusion. The two withdrawn AUC
differences are not decomposed further here, because they differ from the present result in more than
the inputs.

% R3-5 (REPO/docs/tog_manuscript/reviewer_response_matrix.md): TreeSHAP kept only as leak forensics.
% Filled 2026-09-14 by item F2-fill-and-gloss (was \pending{shap_forensics}): REV/A7-shap-forensics/rerun.json, the r3 re-run
%   (REV/queue_state/A7-shap-leak-forensics-v33-m20000-r3.ok 16:06; .done rc 0, final true), script
%   REPO/scripts/run_shap_leak_forensics_v33.py sha1 a9d1362 = committed version (commit 8978bfb).  Its provenance lists dirty
%   manuscript files and scripts/run_leak_ablation.py, which the script does not load (it loads build_corpus_shard,
%   run_deep_tabular_baselines and run_scale_decomposition).  The earlier shap_leak_forensics_v33_m20000.json (05:29, from a
%   mid-edit tree) gave identical values: item F3-fill-results (2026-09-14) compared every JSON leaf of the two files; the 57
%   that differ are timings, argv, output paths and git provenance, and every key cited below is equal to the last digit.
%   rerun.json is the cited file.  Keys: sample.n_matches_sampled 20000; split test rows 15911,
%   matches 5695 (patch 15.14 / 15.15 / 15.16, early stopping on 15.15); test_auc.overall auc_clean 0.647144, auc_leaky
%   0.663690, diff +0.016547, ci_diff [+0.012196, +0.021074] (2,000 match bootstraps); shap.n_rows 10000;
%   leak_share.union share_leaky 0.091930 [0.090298, 0.093523], share_clean 0.004567 [0.004463, 0.004672], diff +0.087363
%   ci_diff [+0.085716, +0.088987] (2,000 bootstraps over the 4,912 explained matches); leak_share.TIME_NORM_ABSOLUTE
%   share_leaky 0.089712, share_clean 0.001322.  Printed as 9.2 %, 0.5 %, 8.7 pp [8.6, 8.9], 9.0 %.
%   Method: docstring "Why SHAP, and only here" (Shapley value definition; mean |SHAP| as global importance; TreeSHAP of
%   Lundberg et al. 2020 via Booster.predict(pred_contrib=True), verified against the definition).
\emph{Where the leak sat in the model.} The ablation says how much the leaking inputs were worth, not
which columns the model used. Attribution answers the second question. A SHAP value splits one prediction
of a model into additive contributions, one per input column, using the Shapley value from cooperative
game theory; for tree models it can be computed exactly~\cite{lundberg2020local}. Averaging the absolute
contributions over many engagements shows which columns a model relies on. We use attribution only for
this forensic purpose: the same engagements are built with and without the leaking inputs, a model is
fitted to each build, and the share of attribution that falls on the columns the leaks change is compared.
On 20,000 sampled matches under the patch holdout, the model given both leaking inputs places 9.2\% of
its attribution on those columns and the audit-fixed model 0.5\%, a paired difference of 8.7 percentage
points (95\% interval 8.6 to 8.9); 9.0 of the 9.2\% fall on game time divided by match length
(Section~\ref{sec:limit-shap}). We do not use attributions to explain what decides engagements.

% -----------------------------------------------------------------------------
\subsection{Class Comparisons Depend on the Teamfight Cut}\label{sec:pred-cut}

% Diagonal and marginal: REV/A6-definition-sensitivity/scale_participation_v33.json populations.labelled.
%   participation: troughs.diagonal peaks [2, 5], primary.k 4, depth 0.146567 (= 1 - 32835 / 38474),
%   bootstrap.primary_location_share {"4": 1.0}, share_4v4_below_3v3_and_5v5 1.0, primary_depth ci_2.5
%   0.134287, ci_97.5 0.160572 (n_boot 1000, match-clustered); diagonal_share_of_population 2v2 0.14698,
%   3v3 0.076145, 4v4 0.061657, 5v5 0.072245 (also printed in sec_definition.tex).
%   troughs.smaller_side peaks [2], troughs [], bootstrap.primary_location_share {"none": 1.0}.
% Class shares at the cuts (stored count -1 read as zero participants, so the 42 rows are picks;
%   scale_participation_v33.json definitions.negative_counts): cuts."3"/"4"/"5":
%   pick n 101840 share 0.191232; skirmish 214295 / 320878 / 392233, share 0.402396 / 0.602535 / 0.736523;
%   teamfight 216412 / 109829 / 38474, share 0.406372 / 0.206233 / 0.072245; largest distance from a share to
%   an end of its share_ci is 0.134 percentage points (computed 2026-09-14 over the nine class-cut cells).
Table~\ref{tab:pred-headline} places the teamfight class at $n_{\min} \ge 4$. The data do not fix that
cut. Among engagements in which both sides field the same number of champions, 2v2 is the most common
(14.7\% of all engagements); between 3v3 (7.6\%) and 5v5 (7.2\%, a second mode at the edge of the range),
4v4 (6.2\%) is a trough. The trough lies at 4v4 in all of 1,000 match-bootstrap resamples (depth 0.147,
95\% interval 0.134 to 0.161, where depth is one minus the 4v4 count divided by the smaller of the two mode
counts), so the equal-sides counts support two kinds of multi-champion engagement. The classes, however,
are defined on $n_{\min}$, the smaller side's count, whose distribution has a single peak at two and no
trough in any resample; and a boundary placed at the 4v4 trough could put 4v4 on either side of it. The
data therefore single out none of the cuts at three, four and five participants, and four is a convention
(Section~\ref{sec:definition}). The cut moves many engagements: the teamfight share is 40.6\%, 20.6\% and
7.2\% at the three cuts and the skirmish share 40.2\%, 60.3\% and 73.7\%, while picks stay at 19.1\%
(each share's 95\% bootstrap interval extends at most 0.14 percentage points on either side).

% Levels: REV/A6-definition-sensitivity/scale_cut_sensitivity_v33.json participation.levels."1".."5":
%   n 101798 / 214295 / 106583 / 71355 / 38474; share 0.191153 / 0.402396 / 0.200138 / 0.133988 / 0.072245;
%   auc.point 0.678869 / 0.667989 / 0.654328 / 0.676101 / 0.689462; ci_2.5 0.675723 / 0.665453 / 0.651202 /
%   0.672395 / 0.684322; ci_97.5 0.682134 / 0.670122 / 0.657789 / 0.679746 / 0.694971 (n_boot 400).
%   levels."0": the 42 rows whose stored count is -1 (a count of zero, see the headline comment), no AUC.
% Gaps: participation.cuts."3"/"4"/"5".gaps:
%   pick_minus_teamfight point +0.010759 / -0.001925 / -0.010499; ci_2.5 +0.006976 / -0.006469 / -0.016997;
%     ci_97.5 +0.014719 / +0.002816 / -0.004128.
%   pick_minus_skirmish point +0.010974 / +0.015636 / +0.013280; ci_2.5 +0.007275 / +0.012308 / +0.009865;
%     ci_97.5 +0.014854 / +0.019545 / +0.017094.
%   skirmish_minus_teamfight point -0.000215 / -0.017561 / -0.023779; ci_2.5 -0.003547 / -0.020804 /
%     -0.029328; ci_97.5 +0.002947 / -0.013618 / -0.018489.
%   reproduces_main_loop_table.by_reading.zero.*.reproduced_to_4dp true (sec_intro.tex, sec_definition.tex).
%   Independent point check in sec_definition.tex ("% cmd" line): {3: 0.0108, 4: -0.0019, 5: -0.0105}.
\begin{table*}[t]
\centering
\caption{Scale classes and the teamfight cut, five-fold cross-validation grouped by match (the predictions
of Table~\ref{tab:pred-headline}, re-scored). Left: AUC at each value of $n_{\min}$, the smaller side's
number of participating champions. Right: paired differences between class AUCs when the teamfight class
begins at $n_{\min} = c$; picks are $n_{\min} \le 1$ and skirmishes $2 \le n_{\min} < c$. The 42
engagements whose stored participation count is $-1$ (a side with no participating champion) are in no
row on the left and count as picks on the right. Match bootstrap, 400 resamples.}
\label{tab:pred-cut}
\scriptsize
\setlength{\tabcolsep}{4pt}
\begin{tabular}{@{}crrc@{\hspace{1.5em}}cccc@{}}
\toprule
\multicolumn{4}{@{}c}{AUC by smaller-side participants} & \multicolumn{4}{c@{}}{Differences between class AUCs [95\% CI]} \\
\cmidrule(r{0.75em}){1-4}\cmidrule(l){5-8}
$n_{\min}$ & $n$ & Share (\%) & AUC [95\% CI] & Cut $c$ & Pick $-$ teamfight & Pick $-$ skirmish & Skirmish $-$ teamfight \\
\midrule
1 & 101,798 & 19.12 & 0.6789 [0.6757, 0.6821] & 3 & $+0.0108$ [$+0.0070$, $+0.0147$] & $+0.0110$ [$+0.0073$, $+0.0149$] & $-0.0002$ [$-0.0035$, $+0.0029$] \\
2 & 214,295 & 40.24 & 0.6680 [0.6655, 0.6701] & 4 & $-0.0019$ [$-0.0065$, $+0.0028$] & $+0.0156$ [$+0.0123$, $+0.0195$] & $-0.0176$ [$-0.0208$, $-0.0136$] \\
3 & 106,583 & 20.01 & 0.6543 [0.6512, 0.6578] & 5 & $-0.0105$ [$-0.0170$, $-0.0041$] & $+0.0133$ [$+0.0099$, $+0.0171$] & $-0.0238$ [$-0.0293$, $-0.0185$] \\
4 & 71,355 & 13.40 & 0.6761 [0.6724, 0.6797] & & & & \\
5 & 38,474 & 7.22 & 0.6895 [0.6843, 0.6950] & & & & \\
\bottomrule
\end{tabular}
\end{table*}

% Non-monotone: level 3 interval [0.6512, 0.6578] overlaps neither level 2 [0.6655, 0.6701] nor level 4
%   [0.6724, 0.6797] (marginal intervals; a conservative reading, no paired level test was run).
Table~\ref{tab:pred-cut} shows why the cut decides the comparison. The AUC does not change monotonically
with $n_{\min}$: it falls from 0.6789 at one participant to 0.6543 at three, whose interval overlaps
neither neighbour's, and rises to 0.6895 at five. A teamfight class that begins at three therefore contains
the hardest level, and one that begins at five contains only the easiest. Accordingly, the
pick-minus-teamfight difference is $+0.0108$ with the cut at three, $-0.0019$ at four and $-0.0105$ at five, and its
interval excludes zero at three and at five, with opposite signs.

% Under a sixth / ninth: max |pick_minus_teamfight point| over cuts 3-5 = 0.010759;
%   0.0677 / 6 = 0.011283 > 0.010759 (ratio 0.010759 / 0.0677 = 0.159);
%   0.097529 / 9 = 0.010837 > 0.010759 (ratio 0.110).  At cut 5: 0.010499 / 0.0677 = 0.155,
%   0.010499 / 0.097529 = 0.108.
The honest summary is therefore not that scale makes no difference. The difference between the pick and
teamfight AUCs moves with a cut that the data do not fix, and it changes sign near the adopted cut. At
every cut its size is under a sixth of the withdrawn draft's $-0.0677$ and under a ninth of the $-0.0975$
of the re-run on the leaking inputs. We make no claim about scale in either direction: neither that larger
engagements are more predictable, nor that predictability does not depend on scale.

% -----------------------------------------------------------------------------
\subsection{What the Scoreboard Explains}\label{sec:pred-scoreboard}

% Readout motivation: sec_intro.tex paragraph 1 ("how much of that likelihood the scoreboard alone implies").
% Scoreboard model: WT/scripts/run_model_comparison_v33.py lead_cols (goldDiff__last, xpDiff__last,
%   avgLevelDiff__last, killDiff_cum__last, csDiff_total__last, aliveDiff__last, time_norm__last,
%   frame_age_s); FEAT/model_comparison_v33_patch.json models.lead_only.n_features 8.  run_model(): lead_only
%   and lgbm_paper share params (n_estimators 400, learning_rate 0.05, num_leaves 31, subsample 0.9,
%   colsample_bytree 0.9, random_state 7) and early_stopping(100) on 15.15; lgbm_deep: n_estimators 3000,
%   learning_rate 0.02, num_leaves 127, min_child_samples 100, reg_lambda 1.0.  time_norm__last under v3.3 =
%   min(t / 45 min, 1).  models.lgbm_paper.n_features = models.lgbm_deep.n_features = 6164.
% Test patch: REV/A10-perclass-contrast/perclass_lead_contrast_v33.json overall n 153816, n_matches 55449,
%   positive_rate 0.504947.
A readout that states which team is ahead is transparent and cheap to build, and the readout proposed in
the Introduction would state how much of the pre-fight likelihood the scoreboard alone implies
(Section~\ref{sec:intro}). We measure this on the patch holdout, with a
\emph{scoreboard model}: the published LightGBM configuration given eight columns, namely the blue-minus-red
differences in gold, experience, average level, kills, creep score and living champions at the last bin,
game time and frame age. It is compared on the same 153,816 test engagements from 55,449 matches with the
same configuration given all 6,164 informative columns, so that only the input differs, and with a
higher-capacity LightGBM (at most 3,000 trees of 127 leaves, learning rate 0.02) on the same columns, the
model behind the figures in the Introduction. Learner capacity itself is examined in
Section~\fwdref{sec:learners}.

\begin{table*}[t]
\centering
\caption{Scoreboard model against full-input LightGBM on the patch holdout (train 15.14, early stopping
15.15, test 15.16), overall and within scale classes at $n_{\min} \ge 4$. Gain is the paired difference in
test AUC from the scoreboard model. Match bootstrap over the 55,449 test matches, 2,000 resamples. Picks,
skirmishes and teamfights are 18.86\%, 60.45\% and 20.67\% of test engagements; the 22 test engagements
whose stored participation count is $-1$ (a side with no participating champion) are in the first row
only.}
\label{tab:pred-scoreboard}
% Source: REV/A10-perclass-contrast/perclass_lead_contrast_v33.json
%   overall.auc.{lead_only,lgbm_paper,lgbm_deep}.point / ci95; overall.contrasts."lead_only->lgbm_paper",
%   "lead_only->lgbm_deep".diff / ci95;
%   by_cut.cut_4.classes.{pick,skirmish,teamfight}.n, share_of_test, auc.*.point / ci95, contrasts.*.diff / ci95.
%   All:       lead 0.626117 [0.623306, 0.628902]; paper 0.666484 [0.663758, 0.669230], gain +0.040367
%              [+0.038501, +0.042187]; deep 0.669212 [0.666560, 0.671929], gain +0.043095 [+0.041159, +0.044949]
%   Pick:      n 29005, share 0.188569; lead 0.601662 [0.595366, 0.608100]; paper 0.676121 [0.669852, 0.682240],
%              gain +0.074458 [+0.069404, +0.079306]; deep 0.679901 [0.673686, 0.685871], gain +0.078239
%              [+0.072900, +0.083322]
%   Skirmish:  n 92988, share 0.604540; lead 0.615592 [0.611832, 0.619385]; paper 0.658946 [0.655363, 0.662291],
%              gain +0.043354 [+0.040730, +0.045945]; deep 0.661880 [0.658315, 0.665331], gain +0.046288
%              [+0.043660, +0.048900]
%   Teamfight: n 31801, share 0.206747; lead 0.671814 [0.665794, 0.677588]; paper 0.679470 [0.673524, 0.685141],
%              gain +0.007656 [+0.004798, +0.010339]; deep 0.680784 [0.674964, 0.686458], gain +0.008970
%              [+0.005968, +0.011841]
%   The overall rows equal REV/A8-input-audit-and-cis/paired_learner_cis_v33.json learners.*.ci95 and
%   comparisons."lead_only->lgbm_paper" / "lead_only->lgbm_deep" to every digit.
% Archived artefact: A10 = REPO/scripts/perclass_lead_contrast_v33.py (commit 0b365d1), final write
%   2026-09-14 15:04:45, seed 7, n_boot 2000, its own seed-8 check "readings that change: none" (validation.
%   monte_carlo_replication max_endpoint_gap_between_seeds 0.000443).  An independent computation by this item
%   (A8 method; not archived) agreed with every class, contrast and gain difference of A10 exactly.
\scriptsize
\setlength{\tabcolsep}{3.5pt}
\begin{tabular}{@{}lrccccc@{}}
\toprule
& & Scoreboard & \multicolumn{2}{c}{Published configuration} & \multicolumn{2}{c@{}}{Higher capacity} \\
& & (8 columns) & \multicolumn{2}{c}{(6,164 columns)} & \multicolumn{2}{c@{}}{(6,164 columns)} \\
\cmidrule(lr){3-3}\cmidrule(lr){4-5}\cmidrule(l){6-7}
Engagements & $n$ & AUC [95\% CI] & AUC [95\% CI] & Gain [95\% CI] & AUC [95\% CI] & Gain [95\% CI] \\
\midrule
All & 153,816 & 0.6261 [0.6233, 0.6289] & 0.6665 [0.6638, 0.6692] & $+0.0404$ [$+0.0385$, $+0.0422$] & 0.6692 [0.6666, 0.6719] & $+0.0431$ [$+0.0412$, $+0.0449$] \\
Pick & 29,005 & 0.6017 [0.5954, 0.6081] & 0.6761 [0.6699, 0.6822] & $+0.0745$ [$+0.0694$, $+0.0793$] & 0.6799 [0.6737, 0.6859] & $+0.0782$ [$+0.0729$, $+0.0833$] \\
Skirmish & 92,988 & 0.6156 [0.6118, 0.6194] & 0.6589 [0.6554, 0.6623] & $+0.0434$ [$+0.0407$, $+0.0459$] & 0.6619 [0.6583, 0.6653] & $+0.0463$ [$+0.0437$, $+0.0489$] \\
Teamfight & 31,801 & 0.6718 [0.6658, 0.6776] & 0.6795 [0.6735, 0.6851] & $+0.0077$ [$+0.0048$, $+0.0103$] & 0.6808 [0.6750, 0.6865] & $+0.0090$ [$+0.0060$, $+0.0118$] \\
\bottomrule
\end{tabular}
\end{table*}

% One model's pick-minus-teamfight AUC difference (paired over classes).  A10 stores the class AUCs and
%   their intervals but not this statistic.
%   Points, from A10 by_cut.cut_c.classes.{pick,teamfight}.auc.<model>.point (durable; differences taken on
%   the unrounded values, so the last digit can differ from the rounded operands):
%     lead_only  c3 0.601662 - 0.654581 = -0.052919;  c4 0.601662 - 0.671814 = -0.070152;
%                c5 0.601662 - 0.683556 = -0.081894
%     lgbm_paper c3 0.676121 - 0.665098 = +0.011022;  c4 0.676121 - 0.679470 = -0.003349;
%                c5 0.676121 - 0.689218 = -0.013098
%   Intervals: the A10 replicate matrix itself, recomputed 2026-09-14 by importing the committed script
%     (REPO/scripts/perclass_lead_contrast_v33.py, commit 0b365d1): load_inputs(default A10 inputs),
%     sets "pick" and "teamfight@c" exactly as main() builds them, B = bootstrap(prep, sets, models, inv,
%     n_groups=55449, n_boot=2000, seed=7), statistic B[:, pick, m] - B[:, teamfight@c, m], interval() =
%     2.5 / 97.5 percentiles.  Check: every class AUC point and ci95 at cuts 3-5 for all three models equals
%     A10 with max abs gap 0.0, so the draws are A10's.  Results (diff [ci95]):
%     lead_only  c3 -0.052919 [-0.061022, -0.045258]; c4 -0.070152 [-0.079034, -0.061360];
%                c5 -0.081894 [-0.092933, -0.070348]
%     lgbm_paper c3 +0.011022 [+0.003366, +0.018449]; c4 -0.003349 [-0.011802, +0.005417];
%                c5 -0.013098 [-0.024167, -0.002030]
%     lgbm_deep  c3 +0.013184 [+0.005813, +0.020817]; c4 -0.000883 [-0.009017, +0.007675];
%                c5 -0.010631 [-0.021547, +0.000468]
%   TODO(orchestrator): add this block to REV/A10-perclass-contrast/perclass_lead_contrast_v33.json (key e.g.
%     by_cut.cut_c.class_auc_differences) and repoint this comment; this item may not write outside its file.
% Teamfight gain of lgbm_paper at cuts 3 and 5: A10 by_cut.cut_3.classes.teamfight.contrasts."lead_only->
%   lgbm_paper" diff +0.010518 ci95 [+0.008132, +0.012703]; cut_5 +0.005662 [+0.002005, +0.009337].
% Gain ordering: REV/A10 .../perclass_lead_contrast_v33.json by_cut.cut_{3,4,5}.ordering.{lgbm_paper,lgbm_deep}
%   .pick_gt_skirmish_gt_teamfight_with_every_pairwise_interval_above_zero true (all six);
%   by_cut.cut_4.gain_differences.lgbm_paper."pick-teamfight" diff +0.066803, ci95 [+0.061050, +0.072514].
% Pick + skirmish shares of the 153,816 test engagements (A10 by_cut.cut_c.classes.*.n):
%   c4 (29005 + 92988) / 153816 = 0.79310; c3 (29005 + 62109) / 153816 = 0.59236;
%   c5 (29005 + 113657) / 153816 = 0.92749.
Table~\ref{tab:pred-scoreboard} gives the result. Over all test engagements the full input adds 0.040 AUC
to the scoreboard model's 0.626, and the addition is uneven. In teamfights the scoreboard model reaches
0.672 and the full input adds 0.008 (95\% interval 0.005 to 0.010). In picks the scoreboard model reaches
only 0.602 and the full input adds 0.074 (0.069 to 0.079); skirmishes lie between, at 0.616 and $+0.043$.
With teamfights at $n_{\min} \ge 4$, the scoreboard model's own AUC is 0.070 lower in picks than in
teamfights (paired interval 0.061 to 0.079), whereas at that cut the published full-input model's pick and
teamfight AUCs cannot be distinguished ($-0.0033$, paired interval $-0.0118$ to $+0.0054$).
% n_min >= 4 qualifier made explicit for the indistinguishability clause (item F2, 2026-09-14): the tie holds only at cut 4;
%   at cut 3 the difference is +0.0110 [+0.0034, +0.0184] and at cut 5 -0.0131 [-0.0242, -0.0020] (next sentence; src above).
For the full-input model this difference depends on the cut, as it does under cross-validation (Section~\ref{sec:pred-cut}): its pick
AUC is above its teamfight AUC by 0.0110 ($+0.0034$ to $+0.0184$) with the cut at three and below it by
0.0131 ($-0.0242$ to $-0.0020$) at five. For the scoreboard model it keeps its sign: the pick AUC stays
below the teamfight AUC by 0.053 at three (0.045 to 0.061) and by 0.082 at five (0.070 to 0.093). In
teamfights the published configuration adds 0.011 to the scoreboard model with the cut at three and 0.006
at five, and the ordering of the gains, largest in picks and smallest in teamfights, holds for both
full-input models at all three cuts, with every pairwise interval above zero.

In plain terms, the rule ``the team that is ahead wins'' ranks the outcomes of teamfights almost as well
as the full model does, and those of picks and skirmishes much worse. Picks and skirmishes are most
engagements wherever the cut is placed: 79.3\% of test-patch engagements with the cut at four, 59.2\% at
three and 92.7\% at five. A readout built on the lead would therefore lose little in large fights and much
in most engagements, and it is in picks and skirmishes that the information beyond the lead matters. This
is a statement about how much of each class's predictable part the lead carries, not a scale gradient of
predictability, on which Section~\ref{sec:pred-cut} makes no claim. Two limits apply. Teamfights happen
later in the match than picks (median cutoff at minute 21.1 against 8.7 in the whole corpus with the cut
at four; Section~\ref{sec:pred-audit}), and we have not separated the effect of game time from that of
the number of participants. And which of the other 6,156 columns carry the gain in picks is not analysed
here.
% 6,156 = 6,164 - 8.

% -----------------------------------------------------------------------------
\subsection{Sensitivity to the Definition's Constants}\label{sec:pred-constants}

% Why: CoG R1 (reviewer_response_matrix.md R1-5b; REPO/scripts/run_gd_sweep_v33.py docstring: "the objection
%   is retired only if the result does not hinge on where they sit").
% Filled 2026-09-14 by item F3-fill-results (was \pending{gxd_sweep} and \pending{presence_gate_v33}).  Every value below was
%   re-read from the two artefacts on 2026-09-14.
% G x D grid: REV/A6-definition-sensitivity/gd_sweep/gd_sweep_summary.json (script REPO/scripts/run_gd_sweep_v33.py; provenance
%   git_commit a5b4965, git_dirty_tracked true; settings_missing []).  provenance: n_matches_sampled 10000 (seed 7, sample_digest
%   51646dbd62f59fa0d9c50186; signature.cache_dir D:/LOL_Project/cache/match_cache_fresh_v3_engage_status13, 210,000 matches per
%   presence_gate_summary.json provenance.n_matches), split patch holdout 15.14 / 15.15 / 15.16, n_boot 1000, tolerance_ms 5000,
%   operating_point G13700_D4264; signature.lgbm = the published configuration (400 trees, learning rate 0.05, 31 leaves,
%   subsample 0.9, colsample_bytree 0.9), early_stopping_rounds 100.  Deviations (file): all 7,106 columns given to LightGBM
%   (train-constant columns are never split on), draws dropped through the label tie policy, one-to-one greedy matching within
%   5 s, per-setting and paired intervals from separate replicate streams (the paired stream resamples the union of both
%   settings' test matches).
%   Operating point settings[G13700_D4264]: counts.n_engagements 27053; counts.by_split.train.rows 9032, test.rows 7298 (2,607
%   matches); test.auc point 0.625191, ci_2.5 0.612380, ci_97.5 0.637923.  Training rows over all 20 settings:
%   counts.by_split.train.rows min 8185 (G18000_D4750), max 10162 (G10000_D3500).
%   Per setting: counts.n_engagements; overlap_engagements.share_of_operating_point_matched; test_rows_vs_operating_point.n_shared
%   and .label_agreement_shared; paired_vs_operating_point.shared_setting_minus_op point [ci_2.5, ci_97.5]; .shared_label_effect.
%   G10000_D3500 eng 30185, op_matched 0.9901, shared 7194, agree 0.9915, shared_setting_minus_op -0.008171 [-0.015477, -0.000819], label -0.004223 [-0.007095, -0.001446]
%   G10000_D4000 eng 29983, op_matched 0.9908, shared 7213, agree 0.9939, shared_setting_minus_op -0.005742 [-0.012303, +0.001264], label -0.003358 [-0.005772, -0.001001]
%   G10000_D4264 eng 29919, op_matched 0.9914, shared 7220, agree 0.9945, shared_setting_minus_op -0.009439 [-0.017364, -0.002415], label -0.003108 [-0.005383, -0.000857]
%   G10000_D4750 eng 29842, op_matched 0.9876, shared 7193, agree 0.9935, shared_setting_minus_op -0.000168 [-0.007524, +0.007046], label -0.002517 [-0.005090, +0.000057]
%   G12000_D3500 eng 28522, op_matched 0.9937, shared 7223, agree 0.9943, shared_setting_minus_op -0.000735 [-0.007638, +0.006331], label -0.001801 [-0.003926, +0.000403]
%   G12000_D4000 eng 28291, op_matched 0.9952, shared 7247, agree 0.9970, shared_setting_minus_op -0.001533 [-0.008439, +0.005813], label -0.001226 [-0.002887, +0.000278]
%   G12000_D4264 eng 28210, op_matched 0.9962, shared 7259, agree 0.9978, shared_setting_minus_op -0.005744 [-0.012442, +0.000713], label -0.000920 [-0.002314, +0.000436]
%   G12000_D4750 eng 28111, op_matched 0.9912, shared 7225, agree 0.9964, shared_setting_minus_op -0.001328 [-0.007613, +0.005118], label -0.000106 [-0.002098, +0.002099]
%   G13700_D3500 eng 27410, op_matched 0.9966, shared 7253, agree 0.9961, shared_setting_minus_op -0.001980 [-0.008548, +0.004345], label -0.001028 [-0.002800, +0.000736]
%   G13700_D4000 eng 27152, op_matched 0.9987, shared 7284, agree 0.9990, shared_setting_minus_op +0.003594 [-0.002659, +0.009402], label -0.000476 [-0.001394, +0.000377]
%   G13700_D4264 eng 27053, op_matched 1.0000, shared 7298, agree 1.0000, shared_setting_minus_op +0.000000 [+0.000000, +0.000000] (operating point)
%   G13700_D4750 eng 26939, op_matched 0.9939, shared 7254, agree 0.9983, shared_setting_minus_op -0.005638 [-0.011571, +0.000046], label +0.000783 [-0.000604, +0.002317]
%   G16000_D3500 eng 26190, op_matched 0.9481, shared 6883, agree 0.9939, shared_setting_minus_op +0.002413 [-0.004257, +0.009603], label -0.000899 [-0.003161, +0.001355]
%   G16000_D4000 eng 25884, op_matched 0.9488, shared 6899, agree 0.9962, shared_setting_minus_op +0.000446 [-0.006787, +0.007394], label -0.000086 [-0.001880, +0.001614]
%   G16000_D4264 eng 25770, op_matched 0.9496, shared 6910, agree 0.9970, shared_setting_minus_op +0.001320 [-0.005739, +0.008552], label +0.000677 [-0.000778, +0.002154]
%   G16000_D4750 eng 25626, op_matched 0.9429, shared 6866, agree 0.9953, shared_setting_minus_op -0.000242 [-0.006893, +0.006800], label +0.001519 [-0.000571, +0.003714]
%   G18000_D3500 eng 25415, op_matched 0.9174, shared 6630, agree 0.9911, shared_setting_minus_op -0.003667 [-0.011821, +0.004158], label -0.000822 [-0.003528, +0.001889]
%   G18000_D4000 eng 25085, op_matched 0.9173, shared 6643, agree 0.9925, shared_setting_minus_op -0.005167 [-0.012050, +0.001931], label -0.000144 [-0.002419, +0.002047]
%   G18000_D4264 eng 24958, op_matched 0.9172, shared 6647, agree 0.9932, shared_setting_minus_op -0.000126 [-0.007218, +0.007421], label +0.000653 [-0.001594, +0.002934]
%   G18000_D4750 eng 24764, op_matched 0.9094, shared 6593, agree 0.9917, shared_setting_minus_op -0.000632 [-0.008250, +0.007018], label +0.001519 [-0.000979, +0.004248]
%   cmd (2026-09-14, over the 19 non-operating settings of settings[*]): n_engagements min 24764 / max 30185; 24764 / 27053 - 1 =
%     -0.0846, 30185 / 27053 - 1 = +0.1158; share_of_operating_point_matched min 0.9094; n_shared min 6593 / max 7284;
%     label_agreement_shared min 0.9911; shared_setting_minus_op min -0.009439 / max +0.003594, 15 of 19 points < 0; intervals
%     entirely below zero only G10000_D3500 and G10000_D4264, none entirely above zero (G13700_D4750 reaches +0.000046, so zero
%     is inside); shared_label_effect intervals entirely below zero only at G10000_D3500 / D4000 / D4264 (-0.004223, -0.003358,
%     -0.003108), none entirely above.  Plateau "10 to 18 s": sec_definition.tex (FB/details_pooled.json temporal.plateau_s).
%   Not printed: per-class test AUCs (test.classes.*; intervals about 0.05 to 0.07 wide at this sample size, e.g. the operating
%     point's teamfight [0.6254, 0.6804]) and full_setting_minus_op (each setting's own test rows: -0.011667 to +0.002621).
% Presence gate: REV/A6-definition-sensitivity/presence_gate/presence_gate_summary.json (script REPO/scripts/
%   run_presence_gate_points.py --engine holdout --summarize; provenance git_commit fa166ec, git_dirty_tracked false; n_matches
%   210000, seed 7, n_boot 1000, tolerance_ms 5000, reference rule_R1600_B15; points_missing []).  M is TF2_MIN_PER_TEAM (script
%   docstring and GATE_KEYS); REQUIRE_ALIVE_PER_TEAM is not changed.  The earlier runs_presence_gate/ comparison used the CoG
%   constants and the leaking inputs and is not cited (sec_definition.tex note).
%   points[rule_R1600_B15]: counts.all.n_engagements 566452; counts.train.n_rows_labelled 187547, counts.test.n_rows_labelled
%     153816; test.auc 0.666484 [0.663722, 0.669230]; reproduces_lgbm_paper: test_auc_reference 0.6664835319909197 =
%     test_auc_this_run, test_auc_diff 0.0, best_iteration 381 = 381, rows_this_run = rows_reference,
%     test_rows_aligned_with_reference true, max_abs_pred_diff 0.0 (reference FEAT/model_comparison_v33_patch.json and .preds.npz).
%   points[ref_R1800_B10]: counts.all.n_engagements 1022489 (1022489 / 566452 = 1.805); counts.train.n_rows_labelled 339470,
%     counts.test.n_rows_labelled 279353; test.auc 0.673588 [0.671787, 0.675575]; paired_vs_reference.full_setting_minus_op
%     +0.007105 [+0.004982, +0.009368]; overlap_engagements.share_of_operating_point_matched 0.913172;
%     test_rows_vs_reference.n_shared 139397, label_agreement_shared 0.991083; auc_vs_reference.setting_shared 0.669312,
%     operating_point_shared 0.665425; paired_vs_reference.shared_setting_minus_op +0.003887 [+0.002550, +0.005307],
%     shared_label_effect -0.001191 [-0.001798, -0.000551], shared_model_effect +0.005078 [+0.003799, +0.006280].
%   points[rule_R1600_B15_M3]: counts.all.n_engagements 76290 (76290 / 566452 = 0.1347); counts.train.n_rows_labelled 25565,
%     counts.test.n_rows_labelled 21173; test.auc 0.678927 [0.672150, 0.686214]; paired_vs_reference.full_setting_minus_op
%     +0.012444 [+0.005689, +0.019138]; test_rows_vs_reference.n_shared 19918, label_agreement_shared 0.997239;
%     auc_vs_reference.setting_shared 0.681144, operating_point_shared 0.687695; paired_vs_reference.shared_setting_minus_op
%     -0.006551 [-0.009725, -0.003525], shared_label_effect -0.000148 [-0.000992, +0.000640], shared_model_effect -0.006402
%     [-0.009395, -0.003522].
%   Deviations (file): shared engagements matched one-to-one on (match id, first-kill time) within 5 s; labels of a shared
%     engagement can differ between points because the label window starts at the cutoff, which moves with B; per-point and
%     paired intervals from separate replicate streams.  Model effect = setting model minus reference model, both scored against
%     the setting's labels (run_gd_sweep_v33.py docstring "shared"), so at ref_R1800_B10 it includes the setting's cutoff
%     (tau = first kill - 10 s) and its training rows.  Not printed: class and presence-class AUCs (test.classes, test.extra).
The result describes the engagements that the definition produces, and the definition has constants. The
kill gap $G$ and the cluster diameter $D$ are estimated from data and the presence radius $R$ and lead $B$
are anchored on game rules (Section~\ref{sec:def-constants}), but a reader may still ask whether the AUC
would move if they sat elsewhere. If it moved much, the result would describe the constants rather than
the task. Two experiments re-detect the engagements with only these constants changed, relabel them, and
score them under the patch holdout with the published LightGBM configuration, so that each setting is
compared with the operating point on the same matches rather than with the cross-validated AUC above.

\emph{Kill gap and diameter.} The first experiment varies $G$ over 10, 12, 13.7, 16 and 18\,s and $D$
over 3,500, 4,000, 4,264 and 4,750\,u, 20 settings, on a random sample of 10,000 of the 210,000 cached
matches drawn with a fixed seed (Table~\ref{tab:pred-gxd}). On so few matches each model is trained on
8,185 to 10,162 engagements instead of 187,547, and the operating point itself reaches a test AUC of only
0.6252 (95\% interval 0.6124 to 0.6379). Absolute AUCs from this sample are therefore not compared with those
elsewhere in this section; only differences from the operating point on the same sample are read. A setting also
changes which engagements exist, so each difference is taken on the test engagements that both the setting
and the operating point detect, paired by match and by the time of the first kill (within 5\,s): the AUC
of the setting's model against the setting's labels minus that of the operating point's model against its
own labels. The two labels of such an engagement agree in at least 99.1\% of cases.

\begin{table}[t]
\centering
\caption{Sensitivity to the kill gap $G$ and the cluster diameter $D$ on a random sample of 10,000 matches,
patch holdout, published LightGBM configuration fitted once per setting. Top: engagements detected. Bottom:
test AUC of the setting minus test AUC of the operating point ($G = 13.7$\,s, $D = 4{,}264$\,u; ---), both
on the test engagements that the two settings share (6,593 to 7,284 of the operating point's 7,298), each
model scored against its own labels. $^\dagger$: the paired 95\% interval (match bootstrap, 1,000
resamples, fitted models held fixed, not adjusted for the 19 comparisons) excludes zero. On this sample the
operating point's test AUC is 0.6252, so only differences are shown.}
\label{tab:pred-gxd}
% Source: REV/A6-definition-sensitivity/gd_sweep/gd_sweep_summary.json settings[*].counts.n_engagements and
%   settings[*].paired_vs_operating_point.shared_setting_minus_op.{point, ci_2.5, ci_97.5}; every value with its
%   interval is listed in the comment block at the start of this subsection.  Dagger = interval entirely below zero
%   (G10000_D3500, G10000_D4264).  6,593 / 7,284 = test_rows_vs_operating_point.n_shared min / max; 7,298 =
%   n_test_operating_point; 0.6252 = settings[G13700_D4264].test.auc.point 0.625191.
\footnotesize
\setlength{\tabcolsep}{5pt}
\begin{tabular}{@{}lrrrr@{}}
\toprule
& \multicolumn{4}{c@{}}{$D$ (u)} \\
\cmidrule(l){2-5}
$G$ (s) & 3,500 & 4,000 & 4,264 & 4,750 \\
\midrule
\multicolumn{5}{@{}l}{\emph{Engagements detected}} \\
10 & 30,185 & 29,983 & 29,919 & 29,842 \\
12 & 28,522 & 28,291 & 28,210 & 28,111 \\
13.7 & 27,410 & 27,152 & 27,053 & 26,939 \\
16 & 26,190 & 25,884 & 25,770 & 25,626 \\
18 & 25,415 & 25,085 & 24,958 & 24,764 \\
\midrule
\multicolumn{5}{@{}l}{\emph{AUC difference from the operating point, shared test engagements}} \\
10 & $-0.0082$\rlap{$^\dagger$} & $-0.0057$ & $-0.0094$\rlap{$^\dagger$} & $-0.0002$ \\
12 & $-0.0007$ & $-0.0015$ & $-0.0057$ & $-0.0013$ \\
13.7 & $-0.0020$ & $+0.0036$ & --- & $-0.0056$ \\
16 & $+0.0024$ & $+0.0004$ & $+0.0013$ & $-0.0002$ \\
18 & $-0.0037$ & $-0.0052$ & $-0.0001$ & $-0.0006$ \\
\bottomrule
\end{tabular}
\end{table}

Over the grid the number of engagements ranges from 24,764 to 30,185 against 27,053 at the operating point,
from $-8.5$\% to $+11.6$\%, and every setting keeps at least 90.9\% of the operating point's engagements.
On the shared test engagements the AUC differences range from $-0.0094$ to $+0.0036$, and 15 of the 19 are
negative. No setting's interval lies above zero. Two lie below it, both at $G = 10$\,s, the lower end of
the 10--18\,s plateau (Section~\ref{sec:definition}): $-0.0082$ ($-0.0155$ to $-0.0008$) at $D = 3{,}500$\,u and
$-0.0094$ ($-0.0174$ to $-0.0024$) at $D = 4{,}264$\,u. At $G = 10$\,s the change of labels costs the most
AUC: scoring the operating point's model against the labels of the settings with $G = 10$\,s and
$D \le 4{,}264$\,u lowers its AUC on the shared engagements by 0.0031 to 0.0042, with intervals below zero.
Labels do not change most often there: at every $D$ the share of shared test engagements whose label
differs from the operating point's is larger at $G = 18$\,s (0.7\% to 0.9\%) than at $G = 10$\,s (0.6\% to
0.8\%).
% Corrected after the F3 fact-check (2026-09-14): the draft said "the labels change most" at G = 10 s, which the artefact
%   does not support.  gd_sweep_summary.json settings[*].test_rows_vs_operating_point.label_agreement_shared:
%   G10000 D3500 / D4000 / D4264 / D4750 0.991521 / 0.993900 / 0.994460 / 0.993466 (label differs on 0.85 / 0.61 / 0.55 /
%   0.65 %); G18000 0.991101 / 0.992473 / 0.993230 / 0.991658 (0.89 / 0.75 / 0.68 / 0.83 %); lowest agreement of all 19
%   settings G18000_D3500.  What is largest at G = 10 s is settings[*].paired_vs_operating_point.shared_label_effect: the four
%   most negative points of the 19 are G10000_D3500 -0.004223 [-0.007095, -0.001446], D4000 -0.003358 [-0.005772, -0.001001],
%   D4264 -0.003108 [-0.005383, -0.000857], D4750 -0.002517 [-0.005090, +0.000057]; next G12000_D3500 -0.001801.
%   Label effect = AUC of the operating point's model against the setting's labels minus against its own labels, on the
%   shared rows (shared_setting_minus_op = shared_label_effect + shared_model_effect, e.g. G10000_D3500 -0.004223 - 0.003948
%   = -0.008171).

The grid shows no more than this. It is one sample of 10,000 matches, and differences between models
trained on the full corpus may be larger or smaller. Each setting is fitted once with one seed, and the
intervals hold the fitted models fixed, so they leave out the variation from refitting. Every difference is
taken against the same fitted model of the operating point, so the 19 differences are not independent: a
favourable fit of that one model would lower all of them together. The intervals are not adjusted for the
19 comparisons. No margin within which a difference counts as negligible was set in advance, so the grid
bounds the differences observed in these ranges and does not establish that the AUC is insensitive to $G$
and $D$; it says nothing about values outside the ranges, about cross-validation or about the scale
classes.

\emph{Presence gate.} The second experiment changes only the presence gate and runs on all 210,000 cached
matches. At the adopted gate ($R = 1{,}600$\,u, $B = 15$\,s, $M = 2$) it reproduces the published
patch-holdout result exactly: the same 153,816 test engagements, the same predictions and a test AUC of
0.6665. At the conference-version gate ($R = 1{,}800$\,u, $B = 10$\,s) the corpus holds 1,022,489
engagements, 1.8 times the 566,452 at the adopted gate, and 91.3\% of the adopted gate's engagements are
among them. The test AUC is 0.6736 (0.6718 to 0.6756) on 279,353 test engagements, $+0.0071$ ($+0.0050$ to
$+0.0094$) above the adopted gate's. On the 139,397 test engagements that both gates accept it is 0.6693
against 0.6654, a paired difference of $+0.0039$ ($+0.0025$ to $+0.0053$). Scoring the adopted gate's model
against the other gate's labels on those engagements changes its AUC by $-0.0012$ ($-0.0018$ to $-0.0006$),
so the gain lies in the model, $+0.0051$ ($+0.0038$ to $+0.0063$). That model differs in two ways that the
run does not separate: it reads the state at a cutoff 10\,s rather than 15\,s before the first kill, and it
is trained on 339,470 engagements rather than 187,547.

Raising $M$ to three leaves 76,290 engagements, 13.5\% of the adopted set. Their test AUC is 0.6789 (0.6721
to 0.6862) on 21,173 test engagements, $+0.0124$ ($+0.0057$ to $+0.0191$) above the adopted gate's. The rise
comes from which engagements remain, not from a better model: on the 19,918 test engagements that both
gates accept, the adopted gate's model reaches 0.6877 and the model trained under $M = 3$, on 25,565
engagements, reaches 0.6811, a paired difference of $-0.0066$ ($-0.0097$ to $-0.0035$), of which the
change of labels accounts for $-0.0001$ ($-0.0010$ to $+0.0006$). As in the grid, each point is fitted
once, the intervals use 1,000 match-bootstrap resamples with the fitted models held fixed, and they are not
adjusted for the comparisons.

The presence gate therefore moves the AUC: by up to $+0.0124$ on each gate's own test engagements, and in
either direction on the engagements that two gates share. The result is not invariant to it. A higher AUC
is not a reason to adopt a gate, because a definition chosen for its predictability would make the
prediction result partly a product of that choice. $R$ and $B$ stay fixed by the game rules and $M$ by the
convention of Section~\ref{sec:definition}, and the AUCs in Tables~\ref{tab:pred-headline},
\ref{tab:pred-cut} and~\ref{tab:pred-scoreboard} are those of that definition. Within the tested
ranges of $G$ and $D$, on the 10,000-match sample, the AUC on shared engagements moved by at most 0.0094.
